\documentclass[11pt]{article}

% Change "review" to "final" to generate the final (sometimes called camera-ready) version.
% Change to "preprint" to generate a non-anonymous version with page numbers.
\usepackage[review,pt]{acl}

% Standard package includes
\usepackage{times}
\usepackage{latexsym}

% For proper rendering and hyphenation of words containing Latin characters (including in bib files)
\usepackage[T1]{fontenc}
% For Vietnamese characters
% \usepackage[T5]{fontenc}
% See https://www.latex-project.org/help/documentation/encguide.pdf for other character sets

% This assumes your files are encoded as UTF8
\usepackage[utf8]{inputenc}

% This is not strictly necessary, and may be commented out,
% but it will improve the layout of the manuscript,
% and will typically save some space.
\usepackage{microtype}

% This is also not strictly necessary, and may be commented out.
% However, it will improve the aesthetics of text in
% the typewriter font.
\usepackage{inconsolata}

%Including images in your LaTeX document requires adding
%additional package(s)
\usepackage{graphicx}

% If the title and author information does not fit in the area allocated, uncomment the following
%
%\setlength\titlebox{<dim>}
%
% and set <dim> to something 5cm or larger.

%\title{Instructions for *ACL Proceedings}
\title{Instruções para artigos em conferências *ACL em português}

% Author information can be set in various styles:
% For several authors from the same institution:
% \author{Author 1 \and ... \and Author n \\
%         Address line \\ ... \\ Address line}
% if the names do not fit well on one line use
%         Author 1 \\ {\bf Author 2} \\ ... \\ {\bf Author n} \\
% For authors from different institutions:
% \author{Author 1 \\ Address line \\  ... \\ Address line
%         \And  ... \And
%         Author n \\ Address line \\ ... \\ Address line}
% To start a separate ``row'' of authors use \AND, as in
% \author{Author 1 \\ Address line \\  ... \\ Address line
%         \AND
%         Author 2 \\ Address line \\ ... \\ Address line \And
%         Author 3 \\ Address line \\ ... \\ Address line}

\author{Primeiro Autor \\
    Instituição / Morada, 1.a linha  \\
    Instituição / Morada, 2.a linha  \\
    Instituição / Morada, 3.a linha  \\
  \texttt{email@domain} \\\And
  Segundo Autor \\
    Instituição / Morada, 1.a linha  \\
    Instituição / Morada, 2.a linha  \\
    Instituição / Morada, 3.a linha  \\
  \texttt{email@domain} \\}

%\author{
%  \textbf{First Author\textsuperscript{1}},
%  \textbf{Second Author\textsuperscript{1,2}},
%  \textbf{Third T. Author\textsuperscript{1}},
%  \textbf{Fourth Author\textsuperscript{1}},
%\\
%  \textbf{Fifth Author\textsuperscript{1,2}},
%  \textbf{Sixth Author\textsuperscript{1}},
%  \textbf{Seventh Author\textsuperscript{1}},
%  \textbf{Eighth Author \textsuperscript{1,2,3,4}},
%\\
%  \textbf{Ninth Author\textsuperscript{1}},
%  \textbf{Tenth Author\textsuperscript{1}},
%  \textbf{Eleventh E. Author\textsuperscript{1,2,3,4,5}},
%  \textbf{Twelfth Author\textsuperscript{1}},
%\\
%  \textbf{Thirteenth Author\textsuperscript{3}},
%  \textbf{Fourteenth F. Author\textsuperscript{2,4}},
%  \textbf{Fifteenth Author\textsuperscript{1}},
%  \textbf{Sixteenth Author\textsuperscript{1}},
%\\
%  \textbf{Seventeenth S. Author\textsuperscript{4,5}},
%  \textbf{Eighteenth Author\textsuperscript{3,4}},
%  \textbf{Nineteenth N. Author\textsuperscript{2,5}},
%  \textbf{Twentieth Author\textsuperscript{1}}
%\\
%\\
%  \textsuperscript{1}Affiliation 1,
%  \textsuperscript{2}Affiliation 2,
%  \textsuperscript{3}Affiliation 3,
%  \textsuperscript{4}Affiliation 4,
%  \textsuperscript{5}Affiliation 5
%\\
%  \small{
%    \textbf{Correspondence:} \href{mailto:email@domain}{email@domain}
%  }
%}

\begin{document}
\maketitle
\begin{abstract}
Este documento suplementa as instruções gerais para autores em conferências sob os auspícios da ACL, contendo instruções específicas para usar as folhas de estilo do \LaTeX{}.
%This document is a supplement to the general instructions for *ACL authors. It contains instructions for using the \LaTeX{} style files for ACL conferences.
O próprio documento segue essas instruções e serve portanto de exemplo para os novos artigos.
%The document itself conforms to its own specifications, and is therefore an example of what your manuscript should look like.
Estas instruções devem ser seguidas tanto para artigos enviados para avaliação como para a versão final de um artigo aceite para publicação.
%These instructions should be used both for papers submitted for review and for final versions of accepted papers.

\end{abstract}

\section{Introdução}
%\section{Introduction}

Estas instruções são para os autores de artigos em \LaTeX~em conferências sob os auspícios da ACL, mas não são suficientes por si só: os autores têm de seguir as instruções gerais para anais ou atas publicadas pela ACL\footnote{\url{http://acl-org.github.io/ACLPUB/formatting.html}}. O presente documento contêm informação adicional no que se refere ao \LaTeX{}
%These instructions are for authors submitting papers to *ACL conferences using \LaTeX. They are not self-contained. All authors must follow the general instructions for *ACL proceedings,\footnote{\url{http://acl-org.github.io/ACLPUB/formatting.html}} and this document contains additional instructions for the \LaTeX{} style files.

Os modelos ("templates") incluem o código \LaTeX{} deste documento (\texttt{acl\_latex.tex}), a folha de estilo para o formatar (\texttt{acl.sty}), um estilo bibliográfico da ACL (\texttt{acl\_natbib.bst}), e a bibliografia para a ACL Anthology (\texttt{anthology.bib}).
%The templates include the \LaTeX{} source of this document (\texttt{acl\_latex.tex}), the \LaTeX{} style file used to format it (\texttt{acl.sty}), an ACL bibliography style (\texttt{acl\_natbib.bst}), an example bibliography (\texttt{custom.bib}), and the bibliography for the ACL Anthology (\texttt{anthology.bib}).

%\section{Engines}
\section{Processadores}

Recomendamos vivamente pdf\LaTeX{} para produzir um ficheiro PDF, em vez de usar \LaTeX{} seguido de dvips+ps2pdf ou dvipdf.
%To produce a PDF file, pdf\LaTeX{} is strongly recommended (over original \LaTeX{} plus dvips+ps2pdf or dvipdf).
A folha de estilo \texttt{acl.sty} também pode ser usada com lua\LaTeX{} e Xe\LaTeX{}, que são especialmente apropriados para artigos com caracteres não latinos.
%The style file \texttt{acl.sty} can also be used with lua\LaTeX{} and Xe\LaTeX{}, which are especially suitable for text in non-Latin scripts.
O arquivo \texttt{acl\_lualatex.tex} neste repositório é um exemplo de como usar \texttt{acl.sty} com lua\LaTeX{} ou Xe\LaTeX{}.
%The file \texttt{acl\_lualatex.tex} in this repository provides an example of how to use \texttt{acl.sty} with either lua\LaTeX{} or Xe\LaTeX{}.

%\section{Preamble}
\section{Preâmbulo}

A primeira linha do ficheiro deve ser
%The first line of the file must be
\begin{quote}
\begin{verbatim}
\documentclass[11pt]{article}
\end{verbatim}
\end{quote}

Para usar o ficheiro de estilo em português na versão de apreciação:
%To load the style file in the review version:
\begin{quote}
\begin{verbatim}
\usepackage[review,pt]{acl}
\end{verbatim}
\end{quote}

Para a versão final, omita a opção \verb|review|:
%For the final version, omit the \verb|review| option:
\begin{quote}
\begin{verbatim}
\usepackage[pt]{acl}
\end{verbatim}
\end{quote}

Para usar Times Roman, coloque o seguinte no preâmbulo:
%To use Times Roman, put the following in the preamble:
\begin{quote}
\begin{verbatim}
\usepackage{times}
\end{verbatim}
\end{quote}
(Também é possível usar txfonts ou newtx.)
%(Alternatives like txfonts or newtx are also acceptable.)

Veja a fonte \LaTeX{} deste documento para comentários sobre outros pacotes que podem ser úteis.
%Please see the \LaTeX{} source of this document for comments on other packages that may be useful.

Defina o título e o autor usando \verb|\title| e \verb|\author|. Dentro da lista de autores, formate vários autores através de  \verb|\and| ou \verb|\And| ou \verb|\AND|; veja a fonte \LaTeX{} para exemplos.
%Set the title and author using \verb|\title| and \verb|\author|. Within the author list, format multiple authors using \verb|\and| and \verb|\And| and \verb|\AND|; please see the \LaTeX{} source for examples.

A caixa contendo o título e os nomes dos autores é dimensionada para um mínimo de 5 cm, por omissão. Se precisar de mais espaço, inclua o seguinte no preâmbulo: 
%By default, the box containing the title and author names is set to the minimum of 5 cm. If you need more space, include the following in the preamble:
\begin{quote}
\begin{verbatim}
\setlength\titlebox{<dim>}
\end{verbatim}
\end{quote}
em que \verb|<dim>| deve ser substituído por um comprimento, que não deve ser menor do que 5 cm.
%where \verb|<dim>| is replaced with a length. Do not set this length smaller than 5 cm.

\section{O corpo do documento}
%\section{Document Body}
\subsection{Notas de rodapé}
%\subsection{Footnotes}
As notas são inseridas através do comando \verb|\footnote|.\footnote{Esta é uma nota de rodapé.}
%Footnotes are inserted with the \verb|\footnote| command.\footnote{This is a footnote.}

%\subsection{Tables and figures}
\subsection{Tabelas e figuras}

Veja-se a Tabela~\ref{tab:acentos} para um exemplo de uma tabela e sua legenda.
%See Table~\ref{tab:accents} for an example of a table and its caption.
%\textbf{Do not override the default caption sizes.}
\textbf{Não altere os tamanhos das legendas.}
\begin{table}
  \centering
  \begin{tabular}{lc}
    \hline
    \textbf{Comando} & \textbf{Saída} \\
    \hline
    \verb|{\"a}|     & {\"a}           \\
    \verb|{\^e}|     & {\^e}           \\
    \verb|{\`i}|     & {\`i}           \\
    \verb|{\.I}|     & {\.I}           \\
    \verb|{\o}|      & {\o}            \\
    \verb|{\'u}|     & {\'u}           \\
    \verb|{\aa}|     & {\aa}           \\\hline
  \end{tabular}
  \begin{tabular}{lc}
    \hline
    \textbf{Comando} & \textbf{Saída} \\
    \hline
    \verb|{\c c}|    & {\c c}          \\
    \verb|{\u g}|    & {\u g}          \\
    \verb|{\l}|      & {\l}            \\
    \verb|{\~n}|     & {\~n}           \\
    \verb|{\H o}|    & {\H o}          \\
    \verb|{\v r}|    & {\v r}          \\
    \verb|{\ss}|     & {\ss}           \\
    \hline
  \end{tabular}
 % \caption{Example commands for accented characters, to be used in, \emph{e.g.}, Bib\TeX{} entries.}
 \caption{Comandos de exemplo para caracteres acentuados, para serem usados por exemplo em entradas Bib\TeX{}.}
  \label{tab:acentos}
\end{table}

É desejável que as fontes usadas nas figuras sejam as mesmas que as usadas no documento. Veja-se a Figura~\ref{fig:exper} para um exemplo de figura e sua legenda.
%As much as possible, fonts in figures should conform to the document fonts. See Figure~\ref{fig:experiments} for an example of a figure and its caption.

Usando o pacote \verb|graphicx|, pode-se incluir arquivos com figuras dentro do ambiente de figura, num ponto apropriado no texto.
%Using the \verb|graphicx| package graphics files can be included within figure environment at an appropriate point within the text.
O pacote \verb|graphicx| oferece vários argumentos opcionais para controlar a aparência da figura. 
%The \verb|graphicx| package supports various optional arguments to control the appearance of the figure.
É preciso inclui-lo explicitamente  no preâmbulo (depois da declaração \verb|\documentclass| e antes de \verb|\begin{document}|), usando o comando \verb|usepackage{graphicx}|.
%You must include it explicitly in the \LaTeX{} preamble (after the \verb|\documentclass| declaration and before \verb|\begin{document}|) using \verb|usepackage{graphicx}|.

\begin{figure}[t]
  \includegraphics[width=\columnwidth]{example-image-golden}
%  \caption{A figure with a caption that runs for more than one line. Example image is usually available through the \texttt{mwe} package without even mentioning it in the preamble.}
\caption{Uma figura com uma legenda com várias linhas. A imagem de exemplo costuma estar incluída no pacote \texttt{mwe} sem tal ser sequer mencionado no preâmbulo.}
  \label{fig:exper}
\end{figure}

\begin{figure*}[t]
  \includegraphics[width=0.48\linewidth]{example-image-a} \hfill
  \includegraphics[width=0.48\linewidth]{example-image-b}
%  \caption {A minimal working example to demonstrate how to place two images side-by-side.}
\caption{Um exemplo muito simples de como colocar duas imagens ao lado uma da outra.}
\end{figure*}

%\subsection{Hyperlinks}
\subsection{Hiperligações/URL}

Versões anteriores de \LaTeX{} poderiam produzir o seguinte erro de compilação:
%Users of older versions of \LaTeX{} may encounter the following error during compilation:
\begin{quote}
\verb|\pdfendlink| ended up in different nesting level than \verb|\pdfstartlink|.
\end{quote}
Isto acontece quando pdf\LaTeX{} é usado e a citação é atravessada por uma mudança de página. A melhor forma de resolver este problema é atualizar a versão de \LaTeX{} para 2018-12-01 ou posteriores. 
%This happens when pdf\LaTeX{} is used and a citation splits across a page boundary. The best way to fix this is to upgrade \LaTeX{} to 2018-12-01 or later.

%\subsection{Citations}
\subsection{Citações}

\begin{table*}
  \centering
  \begin{tabular}{lll}
    \hline
%    \textbf{Output}           & \textbf{natbib command} & \textbf{ACL only command} \\
   \textbf{Saída}           & \textbf{comando natbib} & \textbf{comando só ACL} \\
    \hline
    \citep{Gusfield:97}       & \verb|\citep|           &                           \\
    \citealp{Gusfield:97}     & \verb|\citealp|         &                           \\
    \citet{Gusfield:97}       & \verb|\citet|           &                           \\
    \citeyearpar{Gusfield:97} & \verb|\citeyearpar|     &                           \\
    \citeposs{Gusfield:97}    &                         & \verb|\citeposs|          \\
    \hline
  \end{tabular}
 % \caption{\label{citation-guide}
 \caption{\label{GuiaCitacao}
%    Citation commands supported by the style file.     The style is based on the natbib package and supports all natbib citation commands. It also supports commands defined in previous ACL style files for compatibility.
 Comandos para citações presentes na folha de estilos. Os estilos são baseados no pacote natbib, incluindo além destes alguns comandos definidos em eventos anteriores por uma questão de compatibilidade.
  }
\end{table*}

A Tabela~\ref{GuiaCitacao} mostra a sintaxe permitida pelas folhas de estilo.
%Table~\ref{citation-guide} shows the syntax supported by the style files.
Sugerimos o uso dos estilos natbib.
%We encourage you to use the natbib styles.
Para obter citações na forma ``autor (ano)'', use o comando  \verb|\citet| (citar em forma de texto), como citamos aqui o artigo de \citet{Gusfield:97}.
%You can use the command \verb|\citet| (cite in text) to get ``author (year)'' citations, like this citation to a paper by \citet{Gusfield:97}.
Use o comando \verb|\citep| (citar dentro de parênteses) para obter citações na forma ``(autor, ano)'' \citep{Gusfield:97}.
%You can use the command \verb|\citep| (cite in parentheses) to get ``(author, year)'' citations \citep{Gusfield:97}.
Finalmente, para obter citações na forma ``autor, ano'', úteis quando se faz citações no contexto de um parêntesis, use o comando \verb|\citealp| (por exemplo \citealp{Gusfield:97}).
%You can use the command \verb|\citealp| (alternative cite without parentheses) to get ``author, year'' citations, which is useful for using citations within parentheses (e.g. \citealp{Gusfield:97}).

%A possessive citation can be made with the command \verb|\citeposs|.
%This is not a standard natbib command, so it is generally not compatible with other style files.

%\subsection{References}
\subsection{Referências}

\nocite{Ando2005,andrew2007scalable,rasooli-tetrault-2015}

As folhas de estilo do \LaTeX{} e do Bib\TeX{} que fornecemos seguem em geral o formato bibliográfico da APA (American Psychological Association).
%The \LaTeX{} and Bib\TeX{} style files provided roughly follow the American Psychological Association format.
Se o seu ficheiro bib for chamado \texttt{custom.bib},
colocar as seguintes instruções antes de quaisquer apêndices gerará a secção de referências:
%If your own bib file is named \texttt{custom.bib}, then placing the following before any appendices in your \LaTeX{} file will generate the references section for you:
\begin{quote}
\begin{verbatim}
\bibliography{custom}
\end{verbatim}
\end{quote}

È possível obter a ACL Antology completa como um arquivo Bib\TeX{} de \url{https://aclweb.org/anthology/anthology.bib.gz}.
%You can obtain the complete ACL Anthology as a Bib\TeX{} file from \url{https://aclweb.org/anthology/anthology.bib.gz}.
Em alternativa, pode incluir a antologia e o seu próprio ficheiro .bib, usando os seguintes comandos: 
%To include both the Anthology and your own .bib file, use the following instead of the above.
\begin{quote}
\begin{verbatim}
\bibliography{anthology,custom}
\end{verbatim}
\end{quote}

Consulte a secção~\ref{sec:bibtex} para preparar arquivos em Bib\TeX{}.
%Please see Section~\ref{sec:bibtex} for information on preparing Bib\TeX{} files.

%\subsection{Equations}
\subsection{Equações}

Veja um exemplo de equação:
%An example equation is shown below:
\begin{equation}
  \label{eq:exemplo}
  A = \pi r^2
\end{equation}

As etiquetas para numerar equações, secções, subsecções, figuras e tabelas são definidas com o comando \verb|\label{etiqueta}|. Para as referir no texto, usa-se o comando \verb|\ref{etiqueta}|.
%Labels for equation numbers, sections, subsections, figures and tables are all defined with the \verb|\label{label}| command and cross references to them are made with the \verb|\ref{label}| command.

Este é um exemplo que refere a Equação~\ref{eq:exemplo}.
%This an example cross-reference to Equation~\ref{eq:example}.

%\subsection{Appendices}
\subsection{Apêndices}

Use \verb|\appendix| antes de cada secção de apêndice para mudar a numeração para letras. Veja-se Apêndice~\ref{sec:apendice} para um exemplo.
%Use \verb|\appendix| before any appendix section to switch the section numbering over to letters. See Appendix~\ref{sec:appendix} for an example.

%\section{Bib\TeX{} Files}
\section{Arquivos Bib\TeX{}}
\label{sec:bibtex}

Dados que as entradas em Bib\TeX{} não podem usar Unicode, e que além disso algumas formas de codificar caracteres especiais atrapalham a ordem alfabética no Bib\TeX{}, recomendamos que usem a forma de codificar caracteres especiais apresentada na Tabela~\ref{tab:acentos}.
%Unicode cannot be used in Bib\TeX{} entries, and some ways of typing special characters can disrupt Bib\TeX's alphabetization. The recommended way of typing special characters is shown in Table~\ref{tab:accents}.

Inclua sempre que possível DOI ou URL nas entradas em Bib\TeX{}.
%Please ensure that Bib\TeX{} records contain DOIs or URLs when possible, and for all the ACL materials that you reference.
Use o campo \verb|doi| para DOIs e \verb|url| para URLs.
%Use the \verb|doi| field for DOIs and the \verb|url| field for URLs.
Se uma entrada em Bib\TeX{} tiver um campo DOI ou URL, o título do artigo apontará para o próprio artigo, usando o pacote de \LaTeX{} {\tt hyperref}.
%If a Bib\TeX{} entry has a URL or DOI field, the paper title in the references section will appear as a hyperlink to the paper, using the hyperref \LaTeX{} package.

%\section*{Limitations}
\section*{Limitações}

Este documento não inclui os requisitos sobre conteúdo do artigo para nenhuma conferência sob os auspícios da ACL. Verifique as instruções para autores no que respeita ao tamanho de páginas máximo, a necessidade de incluir uma secção de ``Limitações'', etc.
%This document does not cover the content requirements for ACL or any other specific venue.  Check the author instructions for information on maximum page lengths, the required ``Limitations'' section, and so on.

%\section*{Acknowledgments}
\section{Agradecimentos}
A versão inglesa deste documento foi adaptada por Steven Bethard, Ryan Cotterell e Rui Yan das instruções de conferências ACL anteriores, nomeadamente a ACL 2029, por Douwe Kiela e Ivan Vuli\'{c}, a NAACL 2019 por Stephanie Lukin e Alla Roskovskaya, a ACL 2018 por Shay Cohen, Kevin Gimpel e Wei Lu, a NAACL 2018 por Margaret Mitchell e Stephanie Lukin, as sugestões para 
Bib\TeX{} na (NA)ACL 2017/2018 por Jason Eisner, a
ACL 2017 por Dan Gildea e Min-Yen Kan,
a NAACL 2017 por Margaret Mitchell,
a ACL 2012 por Maggie Li e Michael White,
a ACL 2010 por Jing-Shin Chang e Philipp Koehn,
a ACL 2008 por Johanna D. Moore, Simone Teufel, James Allan e Sadaoki Furui, a ACL 2005 por Hwee Tou Ng e Kemal Oflazer, a ACL 2002 por Eugene Charniak e Dekang Lin, e os formatos de conferências ACL e EACL anteriores desenvolvidos por várias pessoas, incluindo John Chen, Henry S. Thompson e Donald Walker.

%This document has been adapted by Steven Bethard, Ryan Cotterell and Rui Yan from the instructions for earlier ACL and NAACL proceedings, including those for ACL 2019 by Douwe Kiela and Ivan Vuli\'{c}, NAACL 2019 by Stephanie Lukin and Alla Roskovskaya, ACL 2018 by Shay Cohen, Kevin Gimpel, and Wei Lu, NAACL 2018 by Margaret Mitchell and Stephanie Lukin, Bib\TeX{} suggestions for (NA)ACL 2017/2018 from Jason Eisner, ACL 2017 by Dan Gildea and Min-Yen Kan, NAACL 2017 by Margaret Mitchell, ACL 2012 by Maggie Li and Michael White, ACL 2010 by Jing-Shin Chang and Philipp Koehn, ACL 2008 by Johanna D. Moore, Simone Teufel, James Allan, and Sadaoki Furui, ACL 2005 by Hwee Tou Ng and Kemal Oflazer, ACL 2002 by Eugene Charniak and Dekang Lin, and earlier ACL and EACL formats written by several people, including John Chen, Henry S. Thompson and Donald Walker.

Elementos adicionais provêm das instruções da \emph{International Joint Conference on Artificial Intelligence} e da \emph{Conference on Computer Vision and Pattern Recognition}.
%Additional elements were taken from the formatting instructions of the \emph{International Joint Conference on Artificial Intelligence} and the \emph{Conference on Computer Vision and Pattern Recognition}.

A conversão para a língua portuguesa foi feita por Eugénio Ribeiro, e a tradução do presente texto por Diana Santos.
% Bibliography entries for the entire Anthology, followed by custom entries
%\bibliography{anthology,custom}
% Custom bibliography entries only
%\bibliographystyle{fullname} 

\bibliography{custom}

\appendix

%\section{Example Appendix}
\section{Exemplo de apêndice}
%\label{sec:appendix}
\label{sec:apendice}

%This is an appendix.
Isto é um apêndice.

\end{document}
